\section{Key findings}
\label{sec:conclusion-findings}

\begin{enumerate}
    \item \textbf{LP-M\textsuperscript{3}N is empirically tight
    on chains.} On the linear-chain HMC benchmark the
    LP-relaxed loss converges to a test Hamming within $0.003$
    of exact structured-hinge training and within $0.013$ of
    the Forward-Backward Bayes-optimal floor, for both the
    linear and the LeNet-5 backbones.

    \item \textbf{The M\textsuperscript{3}N predictor recovers
    the all-different structure from data.} On symbolic Sudoku,
    the learned pairwise matrix matches the all-different shape
    (zero on the diagonal, positive off-diagonal) without any
    hard-coded constraint encoding, and the per-puzzle $0/1$
    error reaches zero from $N = 100$ onwards.

    \item \textbf{Joint structured training beats two-stage
    OCR + solver.} At matched training-set sizes, the
    M\textsuperscript{3}N predictor reduces the per-puzzle
    $0/1$ error from $0.255$ to $0.051$ at $N = 5000$ ---
    three to five times below the OCR baseline across all
    $N$ values tested.

    \item \textbf{Separate learning rates are essential.} A
    backbone-to-pairwise learning-rate ratio of $10\!:\!1$
    resolves a training-dynamics asymmetry that otherwise
    leaves $W$ under-trained and the predictor stuck at a
    per-cell-classifier plateau.
\end{enumerate}
